\section{Background and Related Works}

\subsection{Privacy Loss Distributions (PLD)}

Suppose we have a DP mechanism $\calM$ that outputs a sample from the continuous distribution $P = \calM(D)$ when given database $D$, and outputs a sample from $Q = \calM(D')$ when given $D'$. The $\epsilon$-hockey stick divergence between two distributions $P, Q$ is defined as:
%
\[H_\epsilon(P, Q) = \int_x \max\{P(x) - e^\epsilon Q(x), 0\} \del x\]
\[= \mathbb{E}_{x \sim P}\left[\max\left\{1 -\frac{e^\epsilon}{e^{\ln(P(x)/Q(x))}}, 0\right\}\right] = \mathbb{E}_{x \sim Q}\left[\max\left\{e^{\ln(P(x)/Q(x))} - e^\epsilon, 0\right\}\right].\]
%
A mechanism $\calM$ satisfies $(\epsilon, \delta)$-DP if and only if for all adjacent databases $D, D'$ we have $H_\epsilon(\calM(D), \calM(D')) \leq \delta$. From the definition, we see that to obtain the $\epsilon$-hockey stick divergence between $P$ and $Q$, it suffices to know their \textit{privacy loss distribution} (PLD):

\begin{definition}
The privacy loss random variable for $P$ and $Q$ is given by sampling $x \sim P$, and computing $\ln(P(x)/Q(x))$. The PLD of $P$ and $Q$ is the distribution of this random variable. 
\end{definition}
%
We frequently use the notion of dominating PLDs:
%
\begin{definition}[Definition 7 in \cite{CharacteristicAccounting}]
The PLD of $P, Q$ dominates the PLD of $P', Q'$ if for any $\epsilon, H_\epsilon(P, Q) \geq H_\epsilon(P', Q')$. We will also say random variable $L$ dominates random variable $L'$ if for any $\epsilon, H_\epsilon(L) \geq H_\epsilon(L')$, where $H_\epsilon(L) = \mathbb{E}_{\ell \sim L}\left[\max\left\{1 - e^{\epsilon - \ell}, 0\right\}\right]$.
\end{definition}

Informally, a PLD dominates another PLD if any privacy guarantee satisfied by mechanisms with the dominating PLD is also satisfied by mechanisms with the dominated PLD. In particular, if the PLD of some pair of distributions $P, Q$ dominates the PLDs of all pairs $\calM(D), \calM(D')$ for adjacent $D, D'$, then if $H_\epsilon(P, Q) \leq \delta$, $\calM$ satisfies $(\epsilon, \delta)$-DP. The following lemma shows that composition preserves domination:

\begin{lem}[Theorem 10 in \cite{CharacteristicAccounting}]\label{lem:dominatingcomposition}
Let $\calM_1, \ldots, \calM_k$ be an adaptive sequence of mechanisms, i.e., each mechanism receives the output of all previous mechanism and the database. Suppose for all $i$ and joint outputs $x$ of $\calM_1, \ldots \calM_{i-1}$, the PLD of $\calM_i(x, D)$ and $\calM_i(x, D')$ is dominated by the PLD of $P_i, Q_i$. Then letting $\calM$ be the composition of these mechanisms, the PLD of $\calM(D), \calM(D')$ is dominated by the PLD of $P_1 \times P_2 \times \ldots, Q_1 \times Q_2 \times \ldots$.

Similarly, if $L_1, L_2, \ldots, L_k$ and $L_1', L_2', \ldots, L_k'$ are random variables such that $L_i$ dominates $L_i'$ for all $i$, then $L_1 + L_2 + \ldots + L_k$ dominates $L_1' + L_2' + \ldots + L_k'$.
\end{lem}

In \cite{CharacteristicAccounting}, only the first part of \cref{lem:dominatingcomposition} is stated. However, the proof easily extends to the second part of \cref{lem:dominatingcomposition}.

Finally, the following lemma informally lets us show that domination for the remove adjacency (i.e., $D$ contains an example zeroed out in $D'$) is equivalent to domination for the add adjacency (i.e., $D'$ contains an example zeroed out in $D$). Thus, we usually only need to prove statements under one of the two adjacencies, and it is implied for the other as well.

\begin{lem}[Lemma 29 in \cite{CharacteristicAccounting}]\label{lem:addvsremove}
The PLD of $P, Q$ dominates the PLD of $P, Q'$ if and only if the PLD of $Q, P$ dominates the PLD of $Q', P$.
\end{lem}

\subsection{Privacy Amplification}
Privacy amplification via sampling analyzes the improvement in privacy given by randomly sampling a minibatch of examples instead of choosing it deterministically. Roughly, a $(\epsilon, \delta)$-DP mechanism run on a batch where each example participates with probability $p$ satisfies $(\log(1 - p + p e^\epsilon), \delta)$-DP. The relative improvement from $\epsilon$ to $\log(1 - p + p e^\epsilon)$ gets better as $\epsilon$ gets smaller: $\log(1 - p + p e^\epsilon) \approx p\epsilon$ for $\epsilon < 1$, but $\log(1 - p + p e^\epsilon) \approx \epsilon - \log(1/p)$ for large $\epsilon$. 
The benefits of privacy amplification via sampling in the independent noise setting of DP-SGD, i.e., the decoder matrix $\bfC=\mathbb{I}$, are extremely well-studied \citep{song2013stochastic,BST14,abadi2016deep,mironov2019r,steinke2022composition,koskela2020computing} with tight analyses. In particular, one round of DP-SGD is dominated by the PLD of $N(0, \sigma^2)$ and $(1-p) \cdot N(0, \sigma^2) + p \cdot N(1, \sigma^2)$ and since each round of DP-SGD has independent randomness, composing this PLD with itself $n$ times gives a tight dominating PLD, i.e. tight $(\epsilon, \delta)$ curve, for DP-SGD.